\documentclass[../textbook.tex]{subfiles}
\begin{document}
\bookchapter{数学基础}{ch:rev-probability}

向量承载特征，矩阵实现批量线性变换，梯度描述目标函数对参数的局部变化，概率规定语言模型输出的含义。这些数学对象共同构成语言建模的基础，并用于推导\term{最大似然估计}{Maximum Likelihood Estimation}{MLE}、\term{熵}{Entropy}{}、\term{交叉熵}{Cross-Entropy}{CE}和\term{库尔贝克—莱布勒散度}{Kullback--Leibler Divergence}{KL}。

\begin{assumption}
词表取有限集合；除明确讨论连续情形外，概率均为离散质量。涉及条件概率的条件事件具有正概率；需要有限对数损失时，预测分布在目标支持集上严格为正。涉及统计估计的推导分别给出采样假设。
\end{assumption}

\section{向量、矩阵及线性变换}

\subsection{向量的方向及尺度}
向量可以看作一组有顺序的数，也可以看作空间中的点或方向。设
$x={(x_1,\ldots,x_d)}^{\mathsf T}\in\Real^d$，其欧氏长度和两个向量的内积分别为
\begin{equation}
 \lVert x\rVert_2=\sqrt{\sum_{i=1}^{d}x_i^2},\qquad
 x^{\mathsf T}y=\sum_{i=1}^{d}x_i y_i.
\end{equation}
内积同时包含长度与夹角信息。余弦相似度
$\frac{x^{\mathsf T}y}{\lVert x\rVert_2\lVert y\rVert_2}$
去除了长度影响，但只有在向量非零时才有定义。在语言模型中，一个向量常表示某个词元或位置的特征；坐标本身通常没有独立的固定语义，向量之间的关系才是模型使用的信息。

向量相加要求维数相同。标量乘法改变长度并可能改变方向。逐元素乘法 $x\odot y$ 与内积不同：前者仍是向量，后者将公共索引求和后得到标量。看到一个公式时，先写出输入和输出形状，往往比先代数化简更容易发现错误。

\subsection{矩阵表示批量数据和线性映射}
矩阵 $A\in\Real^{m\times n}$ 有 $m$ 行、$n$ 列。矩阵乘法
\begin{equation}
 y=Ax,\qquad y_i=\sum_{j=1}^{n}A_{ij}x_j
\end{equation}
把 $n$ 维输入映射到 $m$ 维输出。乘法成立的条件是相邻的公共维数相同；公共维被求和，外侧维被保留。若一批样本按行组成 $X\in\Real^{B\times n}$，常将同一变换写为 $Y=XW$，其中 $W\in\Real^{n\times m}$。这只是批处理约定不同，并不改变每个样本执行同一映射的事实。

线性变换满足 $f(ax+by)=af(x)+bf(y)$。神经网络常用的“线性层”实际计算
\begin{equation}
 z=Wx+b,
\end{equation}
含偏置 $b$ 时严格说是\term{仿射变换}{Affine Transformation}{}。连续堆叠多个仿射变换仍可合并为一个仿射变换，因此网络需要激活函数引入非线性。转置 $A^{\mathsf T}$ 交换行列；逆矩阵只对可逆方阵存在。深度学习中通常直接计算线性方程或矩阵乘积，不应为了形式方便显式求逆。

\section{导数、梯度及链式法则}

一元导数描述函数在某一点附近的变化率。多元标量函数 $f:\Real^d\to\Real$ 的梯度为
\begin{equation}
 \nabla_x f={\left(
 \frac{\partial f}{\partial x_1},\ldots,
 \frac{\partial f}{\partial x_d}
 \right)}^{\mathsf T}.
\end{equation}
梯度与输入同形，并指向欧氏空间中函数上升最快的局部方向；负梯度因此是局部下降方向。梯度为零只说明一阶变化消失，不保证该点是全局最小值。

复合函数通过链式法则求导。若 $z=g(x)$、$f=h(z)$，则
\begin{equation}
 \frac{\partial f}{\partial x}
 =\frac{\partial z}{\partial x}^{\mathsf T}
  \frac{\partial f}{\partial z}.
\end{equation}
计算图中的反向传播就是从标量损失出发，沿依赖关系反向应用这条规则。同一变量经多条路径影响损失时，各路径贡献必须相加。矩阵求导、广播和自动微分的具体规则将在\bookxref{ch:rev-tensors}继续说明。

\section{概率基础}

概率描述不确定事件，而统计推断利用有限样本认识其分布。离散样本空间中，概率质量函数满足
\begin{equation}
 p(x)\geq0,\qquad \sum_x p(x)=1.
\end{equation}
条件概率 $p(x\mid y)$ 表示已知 $y$ 后对 $x$ 的重新分配；它与联合概率的关系为
$p(x,y)=p(x\mid y)p(y)$。语言模型正是反复估计“给定已有文本，下一个词元如何分布”。期望是按概率加权的平均，方差衡量随机变量围绕期望的离散程度。后续各节将在这些基本概念上展开最大似然、交叉熵和序列概率。

\section{随机变量及分布}

\subsection{概率质量函数}
\term{随机变量}{Random Variable}{} $X$ 是对不确定结果的数值或符号表示。大写 $X$ 表示随机变量，小写 $x$ 表示其一个可能取值。若 $X$ 取值于有限集合 $\mathcal X$，\term{概率质量函数}{Probability Mass Function}{PMF}写为
\begin{equation}
 p_X(x)=P(X=x),\qquad p_X(x)\ge0,\qquad\sum_{x\in\mathcal X}p_X(x)=1.
\end{equation}
在不会混淆对象时，下标 $X$ 可以省略。模型中的一个词元编号是一次取值，而词表上的概率向量描述所有候选取值的不确定性。不能把某个编号的大小解释为该词元的概率。

例如，若 $\mathcal X=\{a,b,c\}$ 且概率为 $(0.5,0.3,0.2)$，事件 $X\in\{a,c\}$ 的概率是 $0.7$。概率具有可加性，是因为这些离散取值对应互斥事件。两个并非互斥的事件则不能直接将概率相加，否则交集会被重复计算。

本章主要使用有限分布，以便将概念直接连接到词表预测。连续随机变量用\term{概率密度函数}{Probability Density Function}{PDF} $f_X$ 描述区间概率：
\begin{equation}
 P(a\le X\le b)=\int_a^b f_X(x)\dif x,\qquad\int_{-\infty}^{\infty}f_X(x)\dif x=1.
\end{equation}
密度值可以大于一，因为概率由区间积分而非单点密度给出。对具有密度的变量，单点事件的概率为零；这与离散词元概率不同。

\subsection{记号及支持集}
\begin{center}
\begin{tabularx}{\linewidth}{lX}
\toprule 符号 & 含义\\\midrule
$X,Y$；$x,y$ & 随机变量及其取值。\\
$p(x,y)$；$p(y\mid x)$ & 联合概率及给定 $x$ 后的条件概率。\\
$p_*(x)$ & 讨论对象的真实分布，通常未知。\\
$p_\theta(x)$ & 参数为 $\theta$ 的模型分布。\\
$\widehat p_n(x)$ & $n$ 个观测形成的经验频率分布。\\
$\mathbb E_p[g(X)]$ & 函数 $g(X)$ 在分布 $p$ 下的期望。\\
$H(p)$；$H(p,q)$ & 熵；以 $p$ 为目标、$q$ 为预测的交叉熵。\\
$D_{\mathrm{KL}}(p\|q)$ & 从目标分布 $p$ 到分布 $q$ 的 KL 散度。\\
\bottomrule
\end{tabularx}
\end{center}
\term{支持集}{Support}{}是分布赋予正概率的取值集合，记为 $\operatorname{supp}(p)=\{x:p(x)>0\}$。涉及 $\log q(x)$ 的期望若在某个 $p(x)>0$ 的位置遇到 $q(x)=0$，就可能为无穷。这个边界不是公式上的小细节：一个模型若对真实可能出现的结果给出零概率，该结果一旦发生就会产生无穷大的负对数损失。

除特别说明外，本章使用自然对数，信息量单位为 nat。\footnote{改用以二为底的对数时，单位为 bit。两种记法相差常数 $\log 2$；数值比较应保持底数一致。密度的对数还与坐标及计量单位有关，离散熵的全部性质不能不加条件地移到微分熵。}

\section{联合概率、条件概率及链式分解}

\subsection{条件分布}
若 $p_X(x)>0$，条件概率定义为
\begin{equation}\label{eq:rev-conditional}
 p(y\mid x)=\frac{p(x,y)}{p_X(x)},\qquad
 p_X(x)=\sum_y p(x,y).
\end{equation}
条件化把样本空间限制到已知条件，并重新归一化。它不是删除条件变量，也不是默认条件与结果具有因果关系。网页内容与某个答案同时出现，描述统计相关；因果关系还需机制与干预等依据。

将式\eqref{eq:rev-conditional}两边乘以 $p_X(x)$，得到乘法规则
\begin{equation}
 p(x,y)=p_X(x)p(y\mid x).
\end{equation}
交换分解顺序可得贝叶斯公式
\begin{equation}
 p(x\mid y)=\frac{p(y\mid x)p_X(x)}{\sum_{x'}p(y\mid x')p_X(x')},
\end{equation}
要求分母 $p_Y(y)>0$。分子衡量取值 $x$ 与观测 $y$ 的共同可能性，分母汇总全部候选解释，使结果归一化。

\begin{example}[验证器通过后的正确率]
\label{q:ch2:example:01}
设一个验证器判断生成答案是否正确。候选答案的先验正确率为 $0.2$；正确答案有 $0.9$ 的概率通过，错误答案有 $0.1$ 的概率误通过。记正确事件为 $C$，通过事件为 $A$。联合概率为
\begin{align}
 P(C\cap A)&=0.2\times0.9=0.18,\\
 P(C^c\cap A)&=0.8\times0.1=0.08,\\
 P(A)&=0.18+0.08=0.26.
\end{align}
因此
\begin{equation}
 P(C\mid A)=\frac{0.18}{0.26}=\frac9{13}\approx0.6923.
\end{equation}
通过后的正确率高于原来的 $0.2$，却并非 $0.9$。$P(A\mid C)$ 表示正确答案被接受的概率，$P(C\mid A)$ 表示接受结果中正确答案的比例，条件方向不同。基础正确率较低时，错误答案数量较多，即使误通过率较低也会形成不可忽略的污染。

实际系统可从样本中估计这些条件概率；验证器与候选生成器的相关性以及分布变化会影响估计结果。
\end{example}

\subsection{序列联合分布的分解}
对三个随机变量，连续应用乘法规则有
\begin{align}
 p(x_1,x_2,x_3)
 &=p(x_1,x_2)p(x_3\mid x_1,x_2)\\
 &=p(x_1)p(x_2\mid x_1)p(x_3\mid x_1,x_2).
\end{align}
对一般长度 $T$，归纳得到
\begin{equation}\label{eq:rev-prob-chain}
 p(x_{1:T})=\prod_{t=1}^T p(x_t\mid x_{<t}).
\end{equation}
这个分解没有假定各位置独立。相反，右端保留了完整前缀条件。只有进一步假定 $p(x_t\mid x_{<t})=p(x_t)$，才会退化为独立取值乘积。以固定窗口替代完整前缀同样是一项额外限制。

\begin{proposition}[序列分布的归一化]
对于固定长度 $T$ 的离散序列，若每个前缀条件下的词元分布均归一化，则式\eqref{eq:rev-prob-chain}定义的联合分布也归一化。
\end{proposition}
\begin{proof}
从最后一个位置开始求和：
\begin{align}
 \sum_{x_{1:T}}\prod_{t=1}^T p(x_t\mid x_{<t})
 &=\sum_{x_{1:T-1}}\left[\prod_{t=1}^{T-1}p(x_t\mid x_{<t})\right]
 \underbrace{\sum_{x_T}p(x_T\mid x_{<T})}_{1}\\
 &=\cdots=\sum_{x_1}p(x_1)=1.
\end{align}
\end{proof}
这解释了为何逐位置的词表分布能够定义完整序列概率。零概率前缀上的条件分布可以任意补成归一化分布，而不改变联合概率；除法定义只需用于具有正概率的前缀。

\section{期望、方差及经验分布}

\subsection{期望的概率加权}
对有限分布，函数 $g(X)$ 的期望为
\begin{equation}
 \mathbb E_p[g(X)]=\sum_x p(x)g(x).
\end{equation}
期望不要求是变量的一种可能取值。例如公平二元变量取 $0$ 或 $1$，期望为 $0.5$，但一次观测不必出现 $0.5$。对任意常数 $a,b$，期望的线性性给出
$\mathbb E[aU+bV]=a\mathbb E[U]+b\mathbb E[V]$，无须假定 $U,V$ 独立。

记均值为 $\mu=\mathbb E[X]$，方差为
\begin{equation}
 \operatorname{Var}(X)=\mathbb E[(X-\mu)^2]=\mathbb E[X^2]-\mu^2.
\end{equation}
最后一个等式由展开平方并使用 $\mathbb E[X]=\mu$ 得到。若 $U,V$ 相关，则
\begin{equation}
 \operatorname{Var}(U+V)=\operatorname{Var}(U)+\operatorname{Var}(V)
 +2\operatorname{Cov}(U,V),
\end{equation}
其中 $\operatorname{Cov}(U,V)=\mathbb E[(U-\mathbb E U)(V-\mathbb E V)]$。忽略协方差可能严重低估误差，例如同一文档的多个近重复片段并非新的独立证据。

\subsection{有限样本估计}
设 $X_1,\ldots,X_n$ 独立同分布，均值为 $\mu$、有限方差为 $\sigma^2$。样本均值 $\overline X=n^{-1}\sum_iX_i$ 满足
\begin{equation}
 \mathbb E[\overline X]=\mu,\qquad
 \operatorname{Var}(\overline X)=\frac1{n^2}\sum_i\operatorname{Var}(X_i)=\frac{\sigma^2}{n}.
\end{equation}
第一式来自线性性，第二式还使用了不同样本间协方差为零。标准差随 $n^{-\frac{1}{2}}$ 缩小，而不是随 $n^{-1}$ 缩小。将记录复制十次不会产生十倍独立信息。

对离散取值，经验分布为
\begin{equation}
 \widehat p_n(x)=\frac1n\sum_{i=1}^n\mathbf1[X_i=x].
\end{equation}
于是对任意函数 $g$，
\begin{equation}
 \mathbb E_{\widehat p_n}[g(X)]=\sum_x\widehat p_n(x)g(x)
 =\frac1n\sum_{i=1}^n g(X_i).
\end{equation}
经验平均可以严格写成经验分布下的期望。经验分布由有限样本形成，未观察到的取值在总体中仍可能有正概率。后续的训练目标正是由这个可计算的经验平均构成。

\section{最大似然及负对数损失}

\subsection{似然及概率的对象不同}
假设观测 $x^{(1)},\ldots,x^{(n)}$ 在给定参数后独立，模型给数据的似然为
\begin{equation}
 \operatorname{Lik}(\theta;D)=\prod_{i=1}^n p_\theta(x^{(i)}).
\end{equation}
数据固定时，这个表达式被视为参数 $\theta$ 的函数，称为似然。它一般不在参数空间上积分为一，因此不应直接解释为“参数为真的概率”。参数后验还需要先验和相应归一化。

由于对数严格单调，最大化似然与最大化对数似然具有相同的最优解。再乘以负号，得到最小化问题
\begin{equation}
 \theta^*\in\argmin_\theta \mathcal L_n(\theta),\qquad
 \mathcal L_n(\theta)=-\frac1n\sum_{i=1}^n\log p_\theta(x^{(i)}).
\end{equation}
对数将连乘转为求和，既便于分析，也避免大量小概率连乘导致的数值下溢。除以固定的 $n$ 不改变最优解，但改变梯度尺度；优化时仍需与学习率等设置共同考虑。

\subsection{伯努利参数估计}
\begin{example}
\label{q:ch2:example:02}
设 $X\in\{0,1\}$，$p_\theta(X=1)=\theta$，$0<\theta<1$。在 $n$ 次独立观测中出现 $k$ 次一，负对数似然为
\begin{equation}
 J(\theta)=-k\log\theta-(n-k)\log(1-\theta).
\end{equation}
当 $0<k<n$ 时，求导并令导数为零：
\begin{equation}
 J'(\theta)=-\frac{k}{\theta}+\frac{n-k}{1-\theta}=0
 \quad\Longrightarrow\quad \widehat\theta=\frac{k}{n}.
\end{equation}
且 $J''(\theta)=\frac{k}{\theta^2}+\frac{n-k}{(1-\theta)^2}>0$，所以该驻点是唯一最小值。例如十次中六次为一，估计值为 $0.6$。当 $k=0$ 或 $k=n$ 时，若参数域包含端点，最优解分别位于 $0$ 或 $1$；不能在这些情形机械套用内部驻点条件。

端点解也说明经验估计的局限。训练样本没有某个结果时，最大似然可能将其概率设为零。引入先验、平滑或更强的共享结构可以改变这一行为，但这些操作也改变了估计问题，需明确说明。
\end{example}

\subsection{条件监督、回归及多标签}
在分类任务中，通常给定输入 $x_i$，拟合 $p_\theta(y_i\mid x_i)$。二分类令 $p=\sigma(z)=(1+e^{-z})^{-1}$，则伯努利负对数似然为
\begin{align}
 \ell(z,y)&=-y\log p-(1-y)\log(1-p)\\
 &=\log(1+e^z)-yz\\
 &=\max(z,0)-yz+\log(1+e^{-|z|}).\label{eq:rev-stable-bce}
\end{align}
第二行来自代入 Sigmoid，第三行分别对 $z\ge0$ 与 $z<0$ 提取指数中的较大项。最后一式中指数不大于一，避免了很大的正指数。这是损失稳定计算的数学依据。

例如 $z=2$、$y=1$ 时，$p\approx0.8808$，损失为 $\log(1+e^{-2})\approx0.1269$；若标签改为零，损失为 $2+\log(1+e^{-2})\approx2.1269$。同一预测对不同真值应受到不同惩罚。\footnote{对 $y\in[0,1]$ 的软标签，同一式子仍定义两个类别间的交叉熵；此时 $y$ 不再表示一次伯努利取值，而表示目标概率或平均监督。}

对连续目标，若假定 $Y\mid x\sim\mathcal N(f_\theta(x),\sigma^2)$ 且方差固定，则
\begin{equation}
 -\log p_\theta(y\mid x)=\frac12\log(2\pi\sigma^2)
 +\frac{(y-f_\theta(x))^2}{2\sigma^2}.
\end{equation}
去掉与参数无关的项，最小化负对数似然等价于最小化平方误差。若方差也需要学习，第一项便不能删除；离群值和异方差也可能要求不同的分布模型。

互斥多分类用一个总和为一的类别分布，多标签则允许多个标签同时成立。若将各标签建模为\term{条件独立}{Conditional Independence}{}的伯努利变量，其联合似然为各标签似然的乘积，损失为各自二元交叉熵之和。条件独立是所采用的分布结构，并非多标签数据天然满足的性质。

类别不平衡时，若仅把正类损失乘以 $w>0$，总体条件目标变为 $-wp_*\log q-(1-p_*)\log(1-q)$，其中 $p_*=P(Y=1\mid x)$。求导并令其为零，得到最优输出 $q^*=\frac{wp_*}{wp_*+1-p_*}$。除 $w=1$ 或退化端点外，$q^*$ 不再等于原始条件概率。加权目标改变了拟合对象，不能将其输出直接解释为已校准概率；决策阈值也应在独立验证中按错误代价选择。

\section{熵、交叉熵及 KL 散度}

\subsection{熵衡量分布内部的不确定性}
\term{信息熵}{Information Entropy}{}源于对概率不确定性的定量研究\cite{shannon1948communication}。以下用有限分布推导其性质，信息不确定性与文本语义价值须区分。

对概率 $p(x)>0$ 的事件，定义自信息为 $I(x)=-\log p(x)$。较少见的结果具有较大的自信息；两个独立结果联合出现时，其自信息相加，因为 $-\log[p(x)p(y)]=-\log p(x)-\log p(y)$。按真实出现概率平均，得到熵
\begin{equation}
 H(p)=-\sum_xp(x)\log p(x).
\end{equation}
取 $0\log0=0$，这与 $u\log u$ 在 $u\to0^+$ 时的极限一致。离散概率不大于一，故 $H(p)\ge0$；确定性分布的熵为零。

熵描述一个分布自身的不确定性，并不评价模型是否正确。一个模型把错误答案的概率集中到接近一，也可以具有很低的预测熵。低熵表示确信，确信是否合理还需要与目标分布或观测结果比较。

\subsection{交叉熵及散度的分解}
设目标分布为 $p$，模型或编码分布为 $q$，交叉熵定义为
\begin{equation}
 H(p,q)=-\sum_xp(x)\log q(x).
\end{equation}
采样权重来自 $p$，被取对数的预测来自 $q$，两者不可交换。将 $\log q(x)=\log p(x)-\log[\frac{p(x)}{q(x)}]$ 代入，在有限且支持相容时得到
\begin{align}
 H(p,q)&=-\sum_xp(x)\log p(x)+\sum_xp(x)\log\frac{p(x)}{q(x)}\\
 &=H(p)+D_{\mathrm{KL}}(p\|q).\label{eq:rev-cross-entropy-kl}
\end{align}
第一项不随模型 $q$ 改变，第二项衡量预测分布偏离目标分布所带来的额外平均对数损失。因此，对固定目标分布最小化交叉熵，等价于最小化这个方向的 KL 散度。

\begin{proposition}[KL 散度非负性]
对同一有限集合上的概率分布 $p$ 和 $q$，$D_{\mathrm{KL}}(p\|q)\ge0$，且等号成立当且仅当 $p=q$。当 $p(x)>0$ 而 $q(x)=0$ 时，散度取正无穷。
\end{proposition}
\begin{proof}
非负性可以由基本不等式 $-\log u\ge1-u$ 推出。假设在 $p$ 的支持上 $q$ 为正，则
\begin{align}
 D_{\mathrm{KL}}(p\|q)
 &=-\sum_{x:p(x)>0}p(x)\log\frac{q(x)}{p(x)}\\
 &\ge\sum_{x:p(x)>0}[p(x)-q(x)]\\
 &=1-\sum_{x:p(x)>0}q(x)\ge0.
\end{align}
等号成立要求支持上的比值均为一，且 $q$ 没有把额外概率放到支持之外，即 $p=q$。若 $q$ 在 $p$ 的某个正概率位置为零，散度为正无穷，非负性仍成立。\end{proof}
\begin{corollary}[有限集合上的最大熵]
在 $C$ 个取值上，$H(p)\le\log C$：令 $u(x)=\frac{1}{C}$，则 $D_{\mathrm{KL}}(p\|u)=\log C-H(p)\ge0$。最大熵由均匀分布取得。
\end{corollary}

\subsection{KL 散度方向性}
\begin{example}
\label{q:ch2:example:03}
令 $p=(0.75,0.25)$，$q=(0.5,0.5)$，逐项计算得到
\begin{align}
 H(p)&=-0.75\log0.75-0.25\log0.25\approx0.5623,\\
 H(p,q)&=-0.75\log0.5-0.25\log0.5\approx0.6931,\\
 D_{\mathrm{KL}}(p\|q)&=0.75\log1.5+0.25\log0.5\approx0.1308.
\end{align}
因此 $0.6931\approx0.5623+0.1308$，与式\eqref{eq:rev-cross-entropy-kl}一致。反向计算却得到
\begin{equation}
 D_{\mathrm{KL}}(q\|p)=0.5\log\frac{0.5}{0.75}+0.5\log\frac{0.5}{0.25}
 \approx0.1438.
\end{equation}
二者不同，说明 KL 散度一般不对称。它也不满足三角不等式，因而不是通常意义的距离。训练和蒸馏中的 KL 项必须说明哪个分布提供采样权重；改变方向便改变了目标。

\begin{center}
\includegraphics[width=.92\linewidth]{\subfix{../../figures/theory/revision-likelihood-objective.pdf}}
\captionof{figure}{观测数据与模型通过负对数概率形成经验交叉熵。真实分布提供总体目标，有限样本只提供其估计；二者之间的差异不能由训练损失自动消除。}
\end{center}
\end{example}

\section{序列损失}

\subsection{独热目标及 Softmax}
互斥 $C$ 分类的目标类别为 $y$，独热向量为 $e_y$。对预测分布 $q$，
\begin{equation}
 H(e_y,q)=-\sum_{j=1}^C\mathbf1[j=y]\log q_j=-\log q_y.
\end{equation}
因此用整数类别编号索引正确类别，与显式构造独热向量再求交叉熵，在数学上给出同一标量。前者不需要分配整个独热数组。软标签则有多个非零权重，不能再简化为只读取一个类别。

令网络产生实数分数 $z\in\Real^C$，\term{归一化指数函数}{Softmax}{}将其映射到概率单纯形：
\begin{equation}
 q_j=\frac{e^{z_j}}{\sum_ke^{z_k}},\qquad
 \ell(z,y)=-z_y+\log\sum_ke^{z_k}.
\end{equation}
对全部分数减去相同常数不会改变概率。令 $a=\max_kz_k$，有
\begin{equation}
 \log\sum_ke^{z_k}=a+\log\sum_ke^{z_k-a}.
\end{equation}
右端指数都不大于一。稳定实现应计算这种形式，而不是先计算可能溢出的指数再求对数。分数的绝对平移不可识别，真正决定概率的是分数差。具体梯度将在张量与自动微分章推导。

\subsection{条件生成的最大似然}
给定上下文 $c$，长度为 $T$ 的目标 $y_{1:T}$ 满足
\begin{equation}\label{eq:rev-conditional-sequence-mle}
 p_\theta(y_{1:T}\mid c)=\prod_{t=1}^T p_\theta(y_t\mid c,y_{<t}),\qquad
 -\log p_\theta(y_{1:T}\mid c)=-\sum_{t=1}^T\log p_\theta(y_t\mid c,y_{<t}).
\end{equation}
这里先按固定长度讨论；可变长度的完整生成事件还需包含终止词元，其概率处理见\bookxref{ch:rev-autoregressive}。分类是目标仅含一个类别的特例；翻译、问答和指令回答则把类别预测重复到整个目标序列。条件 $c$ 可以进入每一步计算，却不一定属于损失监督区域。仅屏蔽条件位置的损失，并不会删除模型对条件信息的依赖。

\subsection{序列归一化及样本权重}
对文档 $i$ 的有效预测位置集合 $\mathcal T_i$，令 $N_i=|\mathcal T_i|$。序列链式分解给出损失和
\begin{equation}
 S_i=-\sum_{t\in\mathcal T_i}\log p_\theta(x_t^{(i)}\mid x_{<t}^{(i)}).
\end{equation}
总词元平均为 $\mathcal L_{\mathrm{token}}=\frac{\sum_iS_i}{\sum_iN_i}$，文档等权平均则为 $\mathcal L_{\mathrm{doc}}=n^{-1}\sum_i\frac{S_i}{N_i}$，这里要求被平均的文档均有 $N_i>0$。二者只有在特定条件下才相等。

例如两篇文档分别含两个和八个有效目标，平均词元损失分别为一和二，则 $S_1=2$、$S_2=16$。总词元平均是 $\frac{18}{10}=1.8$，文档平均是 $\frac{1+2}{2}=1.5$。后者对短文档的每个词元赋予更大权重。损失归一化并非只影响显示格式，它直接改变不同样本对学习目标的贡献。

填充位置和不参与监督的提示位置应由明确的有效位置集合排除。注意力掩码规定信息可见性，损失掩码规定优化区域，概率公式不能将两者混为一谈。若全部监督位置都被屏蔽，应报告无有效目标，而不能计算除零的平均损失。

\subsection{困惑度的含义及比较条件}
若平均负对数损失使用自然对数，\term{困惑度}{Perplexity}{PPL}定义为
\begin{equation}
 \mathrm{PPL}=\exp(\mathcal L_{\mathrm{token}})
 =\left[\prod_{i,t\in\mathcal T_i}\frac1{p_\theta(x_t^{(i)}\mid x_{<t}^{(i)})}\right]^{\frac{1}{N}},
 \quad N=\sum_iN_i.
\end{equation}
它是正确目标逆概率的几何平均。若模型对每个正确目标都赋予 $\frac{1}{K}$，困惑度为 $K$；一般情况下，它不是每步实际候选数，也不是错误率。指数是单调的，所以在相同目标和归一化下，困惑度与平均负对数损失的排序相同。

跨模型比较时，需要固定语料、分词器、有效位置、上下文和边界处理。分词更粗的模型可能在更少位置上完成同一文本的编码，直接比较每词元损失会改变计量单位。改用每字节或每字符指标也必须明确文本的编码与归一化条件，不能据此忽略输入协议差异。

\chapterreferences
\end{document}
